% Generated draft: requires mathematical and source review.
\documentclass[11pt]{article}
\usepackage[margin=25mm]{geometry}
\usepackage{fontspec}
\setmainfont{Latin Modern Roman}
\usepackage{amsmath,amssymb,mathtools,unicode-math}
\setmathfont{Latin Modern Math}
\usepackage{enumitem,needspace,xurl,hyperref}
\setlength{\parindent}{0pt}
\setlength{\parskip}{5pt}
\setlength{\emergencystretch}{3em}
\newcommand{\N}{\mathbb{N}}
\newcommand{\Z}{\mathbb{Z}}
\newcommand{\Q}{\mathbb{Q}}
\newcommand{\R}{\mathbb{R}}
\newcommand{\C}{\mathbb{C}}
\newcommand{\F}{\mathbb{F}}
\newcommand{\A}{\mathbb{A}}
\newcommand{\PP}{\mathbb{P}}
\newcommand{\E}{\mathbb{E}}
\newcommand{\eps}{\varepsilon}
\newcommand{\dd}{\,\mathrm{d}}
\begin{document}
\section*{UnsolvedMath ID: 8800043}
\textbf{Source problem code:} catalogue entry
\subsection{Definitions and mathematical statement}
Shparlinski asks to prove that the three elliptic-curve sequence families labeled EC-LCG, EC-PG, and EC-NRG in the slides are ``very well distributed''. The excerpt supplies neither their defining formulas nor the quantitative distribution conjecture. The three labels are retained separately; no guessed recurrence is assigned to any of them.

For a finite elliptic-curve group \(E(\F_q)\), distribution of point outputs, coordinate outputs, or extracted bit strings are different questions. More generally, if the specified output sequence takes values \(u_j\) in a finite set \(\Omega\), a counting formulation would compare
\[
 \#\{0\leq j<M:u_j\in A\}
 \quad\text{with}\quad M\frac{|A|}{|\Omega|}
\]
for a specified family of test sets \(A\subseteq\Omega\). This display explains a possible distribution statistic, not the missing claimed bound. Discrepancy, character sums, correlations, and computational pseudorandomness are not interchangeable meanings of ``very well distributed''.

A precise version needs the definitions of all three generators, admissible elliptic curves and parameters, seeds or keys and any averaging over them, the output map, sample lengths relative to periods, test sets, and the intended error estimates or asymptotic regime. Degenerate seeds and short orbits must be dealt with by the original parameter restrictions, not silently included in an unqualified uniform claim. The source's conjecture remains incomplete until those slide definitions and bounds are recovered.
\subsection{Short English statement}
Prove the claimed strong distribution properties for each of EC-LCG, EC-PG, and EC-NRG. Their definitions, allowed parameters, output conventions, and quantitative meaning of well distributed are missing from the extract.
\subsection{Sources}
\begin{itemize}
\item Local UnsolvedMath registry, permanent ID 8800043, source code catalogue entry; attributed to Igor Shparlinski; no proposal year supplied; update 2026-08-06.
\item IPAM, \emph{Workshop on Number Theory and Cryptography --- Open Problems} (2006), attachment 6224, PDF page 36: \url{https://www.ipam.ucla.edu/programs/workshops/workshop-i-number-theory-and-cryptography-open-problems/?tab=schedule}. Registry access date 2026-07-29; PDF and legacy-PPT visual audit. Provenance reports 19 attachments, 17 PDFs and two legacy PPT decks, with repetitions and solved definitions excluded. Present status and redistribution terms require review.
\item The supplied summary credits Lange and Shparlinski with the sequence constructions and mentions discrepancy and correlation results by Lange, Shparlinski, Gutierrez, and Ibeas, without exact papers, formulas, or estimates. Its general open-status assessment is unverified.
\end{itemize}
\end{document}
